% ============================================================
%  Apéndice C — Entrevista de requerimientos con el Dr. Sandino
%  TT 2026-B066 | ESCOM-IPN
% ============================================================

\chapter{Entrevista de requerimientos con el investigador del laboratorio}
\label{ap:entrevista}

El presente apéndice recoge la transcripción de la entrevista de requerimientos
aplicada al Dr.\ César Augusto Sandino Reyes López, investigador del Laboratorio de
Bioquímica Estructural (Sección de Posgrado, ENMyH-IPN). Las respuestas fundamentan los valores numéricos y las decisiones de
diseño documentadas en el capítulo~4 (véase en particular la sección~\ref{sec:umbrales}).

\bigskip
\noindent\textbf{Entrevistado:} Dr.\ César A. Sandino R. L. \\
\textbf{Laboratorio:} Bioquímica Estructural, Sección de Posgrado, ENMyH-IPN

% -----------------------------------------------------------
\section*{Sección 1. Acceso y usuarios del sistema}
% -----------------------------------------------------------

\textbf{P1.} ¿Quién tiene permitido subir videos y consultar resultados en el
laboratorio? ¿Solo usted o también sus estudiantes y tesistas? ¿Preferiría que cada
persona tuviera su propia cuenta o que compartieran una sola cuenta?

\begin{quote}
Todos los que participan en los proyectos podrían subir sus propios videos. Bastaría
con un rol de investigador en general; no se necesitan permisos diferenciados entre
ellos.
\end{quote}

\textbf{P14.} ¿Cuántas personas usarían el sistema al mismo tiempo? ¿Es probable que
dos investigadores suban videos en paralelo o generalmente lo usaría una persona a la vez?

\begin{quote}
Una persona a la vez, seguramente. En estos proyectos trabajan al mismo tiempo máximo
tres personas más el doctor, lo que da cuatro usuarios posibles, pero no de forma
simultánea, sino en el mismo periodo de tiempo.
\end{quote}

% -----------------------------------------------------------
\section*{Sección 2. Protocolo experimental y grabación}
% -----------------------------------------------------------

\textbf{P2.} ¿Con qué dispositivo graban los videos actualmente (celular, tableta,
cámara fija)? ¿La cámara siempre queda en la misma posición entre sesiones o varía?
¿Tienen algún protocolo establecido de iluminación, distancia o encuadre?

\begin{quote}
\textit{No se obtuvo respuesta registrada en la transcripción para esta pregunta.
Pendiente de confirmar en siguiente reunión.}
\end{quote}

\textbf{P3.} ¿Los videos siempre muestran los cuatro cilindros simultáneamente, o hay
casos con menos cilindros o encuadres distintos? ¿Hay experimentos con solo 1, 2 o 3
especímenes por toma?

\begin{quote}
\textit{No se obtuvo respuesta registrada en la transcripción para esta pregunta.
Pendiente de confirmar en siguiente reunión.}
\end{quote}

\textbf{P4.} ¿La duración siempre es exactamente 20 minutos el Día 1 y 5 minutos el
Día 2, o hay variaciones? ¿Existe algún experimento con más de dos sesiones?

\begin{quote}
Sí, siempre son esas duraciones. El Día 1 de 20 minutos sirve para generar estrés en
los especímenes; de esos 20 minutos solo se analizan los primeros 5 del grupo control
para conocer la conducta basal. El video del Día 2 (5 minutos, grabado 24 horas
después) es el que contiene la información del efecto de los tratamientos y es el que
el sistema debe analizar para todos los grupos. No hay experimentos con más de dos
sesiones.
\end{quote}

\textbf{P5.} ¿En qué formato graban actualmente los videos (MP4, MOV, AVI)? ¿Tienen
grabaciones antiguas en otros formatos que necesitarían analizar con el sistema?

\begin{quote}
Siempre MP4. Graban con Mac y siempre guardan como MP4. Todos los videos que tienen
acumulados son MP4.
\end{quote}

% -----------------------------------------------------------
\section*{Sección 3. Conductas y criterios de clasificación}
% -----------------------------------------------------------

\textbf{P6.} Cuando anotan manualmente, ¿usan algún criterio de duración mínima de
un episodio? Por ejemplo, ¿un movimiento de menos de 1 segundo cuenta como
escalamiento o lo ignoran?

\begin{quote}
Sí, la conducta debe mantenerse alrededor de 3 a 5 segundos para contarse como
episodio válido.
\end{quote}

\textbf{P7.} ¿Observan alguna conducta adicional en el FST además de nado,
inmovilidad y escalamiento (sacudidas, buceos, giros)? ¿Necesitarían que el sistema
las registre?

\begin{quote}
No, solo son tres conductas: nado activo, inmovilidad y escalamiento. El buceo podría
ser una variante del nado y se clasificaría como tal. Es la conducta que podría costar
un poco más de trabajo clasificar, pero se trataría como nado.
\end{quote}

\textbf{P8.} Cuando dos evaluadores anotan el mismo video, ¿en qué conducta suelen
diferir más (escalamiento vs.\ nado, o inmovilidad vs.\ nado)? ¿Qué nivel de
concordancia entre el sistema y el experto consideraría aceptable?

\begin{quote}
La mayor diferencia entre analistas ocurre al distinguir escalamiento y nado. Si la
concordancia con el sistema supera el 85\,\% estamos bien, que es más o menos la
variabilidad que hay entre analistas. El objetivo principal no es tanto la velocidad
sino reducir la variabilidad del factor humano.
\end{quote}

% -----------------------------------------------------------
\section*{Sección 4. Grupos experimentales}
% -----------------------------------------------------------

\textbf{PA.} ¿Cómo están estructurados los grupos experimentales? ¿Cuántos
especímenes tiene cada grupo?

\begin{quote}
Siempre se usan entre 6 y 8 especímenes por grupo. Normalmente trabajan con mínimo
tres grupos: (1) grupo control, al que solo se le aplica el nado forzado sin ninguna
intervención farmacológica; (2) grupo de referencia, al que se le administra un
antidepresivo conocido como fluoxetina o desipramina; y (3) grupo de tratamiento
experimental, al que se le aplica la nueva molécula o intervención que se quiere
evaluar (por ejemplo, acupuntura). Los resultados de cada conducta se comparan
estadísticamente entre grupos, obteniendo media, desviación estándar y varianza.
\end{quote}

\textbf{PB.} ¿El grupo control recibe algún placebo o solo entra al cilindro sin
ninguna sustancia?

\begin{quote}
Solo se mete al cilindro, 20 minutos el Día 1 y 5 minutos el Día 2, sin ningún
tratamiento, fármaco ni placebo.
\end{quote}

% -----------------------------------------------------------
\section*{Sección 5. Resultados, reportes y exportación}
% -----------------------------------------------------------

\textbf{P9.} ¿Qué datos necesita ver al final del análisis? ¿Solo el tiempo total por
conducta por animal, o también el porcentaje, número de episodios y duración promedio
de cada episodio?

\begin{quote}
Necesita: (1) el tiempo total de cada conducta por animal y por sesión; (2) el
desglose por minuto de cada conducta, es decir, el porcentaje de nado, escalamiento e
inmovilidad en el minuto 1, 2, 3, 4 y 5. Esto es relevante para la dinámica
farmacológica.
\end{quote}

\textbf{P10.} El archivo CSV que genere el sistema, ¿lo abriría en Excel o lo
importaría a algún software estadístico (SPSS, GraphPad Prism, R)? ¿Necesita alguna
columna o estructura específica para que importe bien?

\begin{quote}
Lo abriría en Excel primero y desde ahí lo importaría a cualquier software
estadístico. La estructura que necesitan es: una cabecera con el tiempo en segundos
(o minutos), y columnas independientes para nado, escalamiento e inmovilidad, cada
una con su valor en segundos.
\end{quote}

\textbf{P11.} Cuando compara Día 1 vs.\ Día 2, ¿lo hace animal por animal o como
promedio del grupo? ¿Necesitaría comparar también entre distintos grupos de
tratamiento dentro del mismo \textit{dashboard}?

\begin{quote}
Analizan estadísticamente por grupo: calculan media y desviación estándar de todas las
ratas del grupo para cada conducta y luego comparan ese grupo contra los demás grupos
de diferentes tratamientos. Normalmente son mínimo tres grupos los que se están
evaluando al mismo tiempo.
\end{quote}

% -----------------------------------------------------------
\section*{Sección 6. Rendimiento, infraestructura y disponibilidad}
% -----------------------------------------------------------

\textbf{P12.} ¿Cuánto tiempo les toma actualmente analizar manualmente un video de
5 minutos con 4 especímenes? ¿Y uno de 20 minutos?

\begin{quote}
El de 20 minutos no lo analizan completo, solo los primeros 5 minutos del grupo
control. Para un video de 5 minutos con 4 ratas, el análisis manual tarda alrededor
de 30 minutos por rata (se repasa el video tres veces, una por conducta). Con 4 ratas,
un solo video toma aproximadamente 2 horas en total. Además cada analista lo hace de
forma independiente, entonces con 3 analistas se multiplica ese tiempo.
\end{quote}

\textbf{P13.} Si el sistema tardara $X$ minutos en procesar un video de 5 minutos
(con 4 especímenes), ¿cuánto sería el máximo aceptable?

\begin{quote}
No tienen mayor problema con el tiempo que tarde el análisis. Lo que más importa es
que haya menos variabilidad, es decir, que el factor del analista influya lo menos
posible. Lo ideal sería que el mismo video analizado múltiples veces arroje siempre
el mismo resultado, con una variabilidad máxima de entre el 5 y el 10\,\%.
\end{quote}

\textbf{P15.} ¿Con qué equipo cuenta el laboratorio para correr el sistema (laptop,
PC de escritorio, servidor)? ¿Alguno tiene GPU? ¿O preferiría que corriera en la nube?

\begin{quote}
De preferencia en la nube, para que los cambios en el equipamiento impacten lo menos
posible. Tienen un par de Mac (una más antigua y una Mac mini más reciente) y un
servidor pequeño (clúster) pendiente de revisar sus especificaciones de tarjeta gráfica
y procesadores. Mejor GPU, mejor rendimiento del clasificador.
\end{quote}

\textbf{P16.} ¿El sistema necesita estar disponible solo en horario laboral o también
fuera de él (por ejemplo, para dejar un análisis largo corriendo durante la noche)?

\begin{quote}
Disponibilidad 24/7 sería ideal, para poder dejar análisis corriendo en la noche.
\end{quote}

% -----------------------------------------------------------
\section*{Sección 7. Datos, retención y volumen}
% -----------------------------------------------------------

\textbf{P17.} ¿Cuánto tiempo necesita conservar los videos originales después del
análisis? El sistema actualmente los borra a los 30 días. ¿Es suficiente ese plazo o
prefiere más tiempo?

\begin{quote}
Ese plazo es suficiente, sí.
\end{quote}

\textbf{P18.} ¿Los videos y resultados son datos sensibles que requieran acceso
restringido entre investigadores del mismo laboratorio, o es aceptable que todos los
usuarios con cuenta puedan ver todos los experimentos?

\begin{quote}
Sería recomendable que todos los que tienen cuenta puedan ver todos los resultados de
todos los experimentos sin restricción entre ellos.
\end{quote}

\textbf{P19.} ¿Cuántos experimentos de FST realizan aproximadamente por mes/semestre?
¿Y cuántos videos acumulados tienen actualmente?

\begin{quote}
No hay una serie mensual programada; depende del avance de las tesis. Por semestre se
realizan entre 3 y 4 experimentos independientes, con promedio de 4 grupos y 32 a 48
especímenes, lo que genera alrededor de 8 videos por experimento. En total: entre 32 y
40 videos nuevos por semestre. De videos acumulados, tienen aproximadamente 80 a 100
videos desde hace 10--12 años, todos ya analizados manualmente, que podrían usarse para
entrenar y probar el sistema.
\end{quote}

% -----------------------------------------------------------
\section*{Sección 8. Preguntas adicionales derivadas de la retroalimentación}
% -----------------------------------------------------------

\textbf{PC.} ¿El sistema debe analizar el video completo de 20 minutos del Día 1 o
solo una porción?

\begin{quote}
No es necesario analizar los 20 minutos completos. Solo se analizan los primeros
5 minutos del grupo control del Día 1, como referencia de conducta basal (espécimen
sin estrés previo, sin intervención). El video principal de análisis para todos los
grupos es siempre el de 5 minutos del Día 2.
\end{quote}

\textbf{PD.} Sobre el entrenamiento del modelo: ¿recomienda usar los reportes manuales
existentes directamente como datos de entrenamiento?

\begin{quote}
No es recomendable usar los reportes completos para entrenamiento porque el modelo
podría sobre-entrenarse hacia esos patrones específicos. Lo ideal es entrenar con clips
cortos donde la conducta sea inequívoca: momentos donde claramente el espécimen está
escalando, claramente nadando o claramente inmóvil. Una vez entrenado, sí se pueden
usar los videos ya analizados manualmente como conjunto de prueba para medir qué tan
cerca llega el modelo (80\,\% o más de concordancia es un umbral razonable).
\end{quote}

\textbf{PE.} ¿Cuál es el nombre correcto del laboratorio colaborador para usarlo en
toda la documentación?

\begin{quote}
Laboratorio de Bioquímica Estructural, Sección de Posgrado, Escuela Nacional de
Medicina y Homeopatía (ENMyH-IPN). El nombre ``Laboratorio de Neurociencia
Conductual'' que aparecía en versiones anteriores es incorrecto.
\end{quote}

\textbf{PF.} ¿Es suficiente referirse solo a ``fluoxetina'' en lugar de
``ISRS~+~fluoxetina'' en la presentación?

\begin{quote}
Sí, es la recomendación. La fluoxetina es el antidepresivo más utilizado en la
clínica. Con decir ``fluoxetina, que es un antidepresivo de referencia'' es suficiente.
Usar la sigla ISRS abre líneas de pregunta sobre farmacología que no son necesarias
para la defensa del TT.
\end{quote}

\textbf{PG.} ¿El Dr.\ Sandino tiene disponibles videos y reportes manuales para
compartir con el equipo como material de entrenamiento?

\begin{quote}
Sí. Ya se subió al canal de Teams la presentación del año anterior con videos
embebidos de las tres conductas. Están localizando los archivos de Excel donde los
alumnos capturaron los análisis manuales para empatar archivo con video. Se comprometió
a compartir esos reportes y también la guía/manual de cómo identificar las conductas.
\end{quote}
